\section{Maps in the Go Programming Language}

Maps are a built-in data structure in the Go programming language
used to store and manage collections of key--value pairs. Each
key is associated with a value, allowing data to be retrieved,
inserted, updated, and deleted efficiently using the corresponding
key. Maps are particularly useful when information must be
organised according to unique identifiers, such as usernames,
product codes, configuration parameters, or database records.\\

Unlike arrays and slices, which organise elements by numerical
index, maps provide access through keys of a specified type.
The key identifies an entry, while the associated value stores
the information linked to that key. Go maps are implemented as
hash tables, although their internal representation and specific
performance characteristics are implementation details.\\

A map provides an abstraction for associative storage and is
commonly used when direct access to values by meaningful keys
is more appropriate than sequential indexing.

\subsection{Map Declaration and Type}

A map type is declared using the \texttt{map} keyword, followed
by the key type in square brackets and the value type.

The general syntax is:

\begin{lstlisting}
map[KeyType]ValueType
\end{lstlisting}

Here, \texttt{KeyType} specifies the type of the keys, while
\texttt{ValueType} specifies the type of the associated values.
All keys in a particular map must have the same type, and all
values must conform to the map's declared value type.\\

For example, a map associating names with ages can be declared
as follows:

\begin{lstlisting}
var ages map[string]int
\end{lstlisting}

In this declaration, the keys are strings and the values are
integers. The variable \texttt{ages} is initially a nil map.
Although its length is zero, it cannot be used to insert
entries until it has been initialised.\\

Map key types must be comparable. Common examples include
integers, strings, Boolean values, pointers, and suitable
struct types. Slices, maps, and functions cannot be used
directly as map keys because they are not comparable types.

For example:

\begin{lstlisting}
var scores map[string]float64
var lookup map[int]string
var flags map[bool]string
\end{lstlisting}

These declarations represent maps with different key and
value types. The types of the keys and values are determined
at compile time.

\subsection{Map Initialisation}

A map must be initialised before entries can be inserted.
Go provides two principal methods for initialising maps:
the built-in \texttt{make} function and map literals.

\subsubsection{Initialisation Using \texttt{make}}

The \texttt{make} function creates an initialised map that
is ready for use.

\begin{lstlisting}
ages := make(map[string]int)
\end{lstlisting}

The resulting map is empty but can immediately accept
key--value assignments. An optional initial capacity hint
can also be supplied:

\begin{lstlisting}
ages := make(map[string]int, 100)
\end{lstlisting}

The second argument provides an indication of the expected
number of entries. It may help the implementation prepare
storage efficiently, but it does not restrict the map to
100 entries.

\subsubsection{Initialisation Using a Map Literal}

A map literal allows entries to be specified directly
during initialisation.

\begin{lstlisting}
ages := map[string]int{
    "Alice": 25,
    "Bob":   30,
    "Carol": 28,
}
\end{lstlisting}

This statement creates a map containing three entries.
The strings are keys, and their associated integers are
the stored values.\\

Map literals are useful when the initial contents are known
in advance, particularly for configuration settings,
lookup tables, and small collections of related data.

\subsection{Inserting and Updating Map Entries}

Entries are inserted into a map by assigning a value to a
key using the indexing operator.

\begin{lstlisting}
scores := make(map[string]int)

scores["Alice"] = 85
scores["Bob"] = 92

fmt.Println(scores["Alice"]) // 85
\end{lstlisting}

If the specified key does not exist, the assignment creates
a new entry. If the key already exists, its associated
value is replaced.

For example:

\begin{lstlisting}
scores["Alice"] = 85
scores["Alice"] = 95

fmt.Println(scores["Alice"]) // 95
\end{lstlisting}

The second assignment updates Alice's score from 85 to 95.
A map cannot contain multiple distinct entries with the
same key; each key identifies at most one current value.\\

Assignment to a nil map is invalid and causes a run-time
panic. Therefore, a map should be initialised before
inserting entries.

\subsection{Accessing Map Values}

A map value can be retrieved by specifying its key in
square brackets.

\begin{lstlisting}
ages := map[string]int{
    "Alice": 25,
    "Bob":   30,
}

fmt.Println(ages["Alice"]) // 25
\end{lstlisting}

If a key is absent, a map lookup returns the zero value
of the map's value type. For example, a lookup in a
\texttt{map[string]int} returns zero when the specified
key does not exist.\\

This behaviour can make it difficult to distinguish an
absent entry from an entry whose value is genuinely zero.
Go addresses this issue through the two-result form of
map indexing.

\subsubsection{Checking Whether a Key Exists}

A map lookup can return both the associated value and
a Boolean indicating whether the key is present.

\begin{lstlisting}
score, exists := scores["Alice"]

if exists {
    fmt.Println("Score:", score)
} else {
    fmt.Println("Entry not found")
}
\end{lstlisting}

The variable \texttt{exists} is \texttt{true} if the key
is present and \texttt{false} otherwise. This approach
allows a program to distinguish an absent key from a
key associated with a zero value.\\

For example:

\begin{lstlisting}
quantities := map[string]int{
    "apples": 0,
}

value, exists := quantities["apples"]
fmt.Println(value, exists) // 0 true

value, exists = quantities["oranges"]
fmt.Println(value, exists) // 0 false
\end{lstlisting}

Both lookups return zero as the value, but their Boolean
results distinguish the two situations.

\subsection{Deleting Map Entries}

Go provides the built-in \texttt{delete} function to
remove an entry from a map.

\begin{lstlisting}
ages := map[string]int{
    "Alice": 25,
    "Bob":   30,
    "Carol": 28,
}

delete(ages, "Bob")

fmt.Println(len(ages)) // 2
\end{lstlisting}

The function accepts a map and a key. If the key exists,
the corresponding entry is removed. If the key is absent,
the operation has no effect. Deleting an entry from a
nil map is also permitted.\\

The built-in \texttt{clear} function can be used to remove
all entries from a map.

\begin{lstlisting}
clear(ages)
fmt.Println(len(ages)) // 0
\end{lstlisting}

The \texttt{clear} operation empties the map without
requiring the variable to be reassigned to a new map.

\subsection{Determining the Size of a Map}

The built-in \texttt{len} function returns the number of
entries currently stored in a map.

\begin{lstlisting}
products := map[string]float64{
    "Laptop":  899.99,
    "Monitor": 249.50,
    "Mouse":    29.99,
}

fmt.Println(len(products)) // 3
\end{lstlisting}

The result represents the number of key--value pairs,
rather than the amount of memory allocated internally.
The length changes when entries are inserted or deleted.\\

For a nil map, \texttt{len} returns zero.

\subsection{Iterating Over Maps}

The \texttt{range} clause can be used to iterate over the
entries of a map. During iteration, it provides each key
and its associated value.

\begin{lstlisting}
scores := map[string]int{
    "Alice": 85,
    "Bob":   92,
    "Carol": 78,
}

for name, score := range scores {
    fmt.Println(name, score)
}
\end{lstlisting}

Each iteration assigns a key to \texttt{name} and its
associated value to \texttt{score}. The loop continues
until the map's entries have been visited.\\

Map iteration order is unspecified. A program must not
assume that entries will be returned in insertion order,
alphabetical order, or any other consistent order.
The order may vary between iterations.\\

If deterministic output is required, the keys can be
collected into a slice, sorted, and then used to retrieve
the associated values.

\begin{lstlisting}
names := make([]string, 0, len(scores))

for name := range scores {
    names = append(names, name)
}

sort.Strings(names)

for _, name := range names {
    fmt.Println(name, scores[name])
}
\end{lstlisting}

This example requires the \texttt{sort} package. Sorting
the keys ensures that the entries are printed in
alphabetical order, independently of the map's iteration
order.

\subsection{Maps and Composite Values}

Map values may be of any type, including arrays, slices,
structs, pointers, interfaces, and other maps. This allows
maps to represent complex data relationships.\\

For example, a map can associate a person's name with
a structure containing several attributes.

\begin{lstlisting}
type Person struct {
    Age    int
    City   string
}

people := map[string]Person{
    "Alice": {
        Age:  25,
        City: "Manchester",
    },
    "Bob": {
        Age:  30,
        City: "London",
    },
}
\end{lstlisting}

A value can be accessed through its key and its fields
can be read directly.

\begin{lstlisting}
person := people["Alice"]
fmt.Println(person.City) // Manchester
\end{lstlisting}

When a map stores struct values, retrieving an entry
produces a copy of the stored struct value. Consequently,
a field of a struct stored directly in a map cannot be
assigned through a map index expression.\\

For example, the following assignment is invalid:

\begin{lstlisting}
// Invalid:
// people["Alice"].Age = 26
\end{lstlisting}

Instead, the value can be retrieved, modified, and assigned
back to the map.

\begin{lstlisting}
person := people["Alice"]
person.Age = 26
people["Alice"] = person
\end{lstlisting}

Alternatively, a map can store pointers to structs, allowing
the pointed-to objects to be modified directly.

\begin{lstlisting}
people := map[string]*Person{
    "Alice": {Age: 25, City: "Manchester"},
}

people["Alice"].Age = 26
\end{lstlisting}

This approach is appropriate when shared mutable objects
are required, although pointer ownership and concurrent
access must be managed carefully.

\subsection{Nested Maps}

A map can contain another map as its value type. This
creates a nested map structure that can represent
hierarchical relationships or multidimensional lookup
tables.\\

For example, a system may store student marks by student
name and subject.

\begin{lstlisting}
marks := map[string]map[string]int{
    "Alice": {
        "Mathematics": 90,
        "English":     85,
    },
    "Bob": {
        "Mathematics": 78,
        "English":     88,
    },
}

fmt.Println(marks["Alice"]["Mathematics"]) // 90
\end{lstlisting}

The outer map associates each student with another map,
which stores subject names and their corresponding marks.\\

When nested maps are created incrementally, each inner
map must be initialised before entries can be inserted
into it.

\begin{lstlisting}
results := make(map[string]map[string]int)

results["Alice"] = make(map[string]int)
results["Alice"]["Mathematics"] = 90
\end{lstlisting}

Without the initialisation of the inner map, assigning
a value through a nil inner map would cause a run-time
panic.

\subsection{Maps and Function Parameters}

Maps can be passed to functions as arguments. A map value
contains a reference to the underlying map data, so copying
the map value does not create an independent copy of all
its entries. Consequently, modifications made through
the parameter can be visible to the caller.

\begin{lstlisting}
func updateScore(scores map[string]int,
    name string, score int) {
    scores[name] = score
}

func main() {
    scores := map[string]int{
        "Alice": 85,
    }

    updateScore(scores, "Alice", 95)

    fmt.Println(scores["Alice"]) // 95
}
\end{lstlisting}

In this example, the function updates the entry in the
map that is also accessible through the caller's variable.
The map value itself is passed by the value, but both values
refer to the same underlying map data.\\

Assigning a new map to the function's local parameter
does not, however, automatically replace the caller's
map variable. If a function needs to create or replace
a map and communicate that new map value to the caller,
it can return the map.

\begin{lstlisting}
func createScores() map[string]int {
    return map[string]int{
        "Alice": 85,
        "Bob":   92,
    }
}
\end{lstlisting}

This pattern is useful for constructing maps within
functions and returning them to other parts of a program.

\subsection{Nil Maps and Empty Maps}

Go distinguishes between nil maps and initialised empty
maps. Both have a length of zero, but their behaviour
differs when entries are inserted.\\

A nil map can be declared without initialisation:

\begin{lstlisting}
var scores map[string]int

fmt.Println(scores == nil) // true
fmt.Println(len(scores))   // 0
\end{lstlisting}

A nil map supports lookups, iterations, deletions, and
length queries. A lookup returns the value type's zero
value when the key is absent. However, assigning a new
entry to a nil map causes a run-time panic.

An empty map can be created using \texttt{make} or an
empty map literal.

\begin{lstlisting}
scores := make(map[string]int)
// Alternatively:
scores = map[string]int{}

scores["Alice"] = 85
\end{lstlisting}

This map is initialised and supports insertion. Choosing
between a nil map and an empty map depends on the
requirements of the program and the conventions of the
surrounding API.

\subsection{Map Equality and Comparability}

Map types are not comparable, except that a map value
can be compared with \texttt{nil}. Therefore, two maps
cannot be compared directly using the equality operator
to determine whether they contain the same entries.

\begin{lstlisting}
first := map[string]int{"Alice": 25}
second := map[string]int{"Alice": 25}

fmt.Println(first == nil) // false

// Invalid:
// fmt.Println(first == second)
\end{lstlisting}

If two maps must be compared by their contents, the
program must compare their lengths, keys, and associated
values explicitly, or use an appropriate comparison
utility where the value types permit it.\\

The inability to compare maps directly reflects their
reference-like behaviour and the absence of a general
language-defined equality operation for map contents.

\subsection{Concurrency and Map Safety}

Built-in Go maps are not safe for concurrent read and
write access without appropriate synchronisation.
If multiple goroutines access a map concurrently and
at least one goroutine modifies it, the program must
coordinate access to prevent data races and invalid
concurrent operations.\\

One common approach is to protect map access with a
mutex from the \texttt{sync} package.

\begin{lstlisting}
type SafeCounter struct {
    mu     sync.Mutex
    counts map[string]int
}

func (c *SafeCounter) Increment(key string) {
    c.mu.Lock()
    defer c.mu.Unlock()

    c.counts[key]++
}
\end{lstlisting}

The mutex ensures that only one goroutine at a time
executes the protected operation. In practical code,
the map must also be initialised before insertion.\\

A complete initialisation might be:

\begin{lstlisting}
counter := SafeCounter{
    counts: make(map[string]int),
}
\end{lstlisting}

For workloads involving specialised concurrent access
patterns, Go also provides the \texttt{sync.Map} type.
However, an ordinary map protected by a mutex is often
simpler and provides clearer control over the relationship
between map operations and other shared state.

\subsection{Performance Considerations}

Maps generally provide efficient insertion, retrieval,
and deletion by key. Their performance is typically
constant-time on average, although individual operations
can take longer depending on collisions, resizing, key
costs, and implementation details.\\

The choice of key type can influence performance.
String keys are useful for descriptive identifiers,
whereas integer keys may be appropriate for numerical
identifiers. Key types must also be comparable, and
equality comparisons for compound key types may involve
multiple fields.\\

When the approximate number of entries is known, providing
an initial capacity \texttt{make} can help reduce
allocation and resizing overhead.

\begin{lstlisting}
users := make(map[int]string, 1000)
\end{lstlisting}

This creates a map intended to hold approximately 1,000
entries. It does not impose a maximum capacity.\\

Maps should be selected when data is best accessed through
keys rather than sequential positions. For workloads that
require ordered traversal, direct indexing, or compact
fixed-size storage, arrays or slices may be more appropriate.

\subsection{Conclusion}

Maps are an essential data structure in Go for representing
associations between unique keys and corresponding values.
They support efficient insertion, retrieval, updating,
deletion, and iteration, making them suitable for lookup
tables, counters, configuration data, caching, and
structured relationships.\\

A thorough understanding of map declaration, initialisation,
key comparability, zero-value lookups, iteration order,
nil-map behaviour, nested structures, and concurrency
is necessary for writing correct and efficient Go programs.
Particular attention should be given to the fact that map
iteration order is unspecified and that built-in maps
require synchronisation when accessed concurrently with
modifications.\\

